\documentclass[11pt]{amsart}
\usepackage{amsmath,amssymb,amsthm}
\usepackage[margin=1.25in]{geometry}

\newtheorem{theorem}{Theorem}
\newtheorem{proposition}[theorem]{Proposition}
\newtheorem{lemma}[theorem]{Lemma}
\theoremstyle{remark}
\newtheorem*{remarks}{Remarks}

\newcommand{\NN}{\mathbb{N}}
\newcommand{\PP}{P}

\title[Dickson's conjecture and sum-sets within the primes]
      {Dickson's conjecture and infinite multidimensional\\ sum-sets within the primes}
\author{David Feldman}
\address{Department of Mathematics and Statistics, University of New
Hampshire, Durham, NH 03824, USA}
\author{Russell Impagliazzo}
\address{Department of Computer Science and Engineering, University of
California San Diego, La Jolla, CA 92093, USA}
\subjclass[2020]{11P32 (11N32, 68V20)}

\begin{document}

\begin{abstract}
Assuming Dickson's conjecture, we show that for every prime $p$ and
every $n\le p-2$ the set $\PP$ of primes contains a sum-set
$\{p\}+E_0+\cdots+E_n$ with each $E_i$ an infinite set of even numbers
containing $0$.  Unconditionally, no such configuration exists once
$n\ge p-1$, so the result is sharp.  In particular the primes contain
no infinite-dimensional sum-set, although random models capturing the
density behavior of the primes do; in this respect the primes fail to
behave like a typical prime-like random set.
\end{abstract}

\maketitle

\section{Introduction}

We shall show that Dickson's conjecture concerning the set $\PP$ of
prime numbers implies that $\PP$ contains certain strikingly highly
structured subsets.  Precise statements follow once we establish
terminology and notation.  Throughout, $\NN$ denotes the natural
numbers including $0$, and $2\NN$ the even natural numbers.

Given a family $\{A_i\}$ of subsets of $\NN$, write $\bigoplus A_i$
for the subset of $\prod A_i$ consisting of vectors with all but
finitely many coordinates equal to $0$.  (For infinite families
$\{A_i\}$ we always assume all but finitely many sets contain $0$.)
The \emph{sum-set} of $\{A_i\}$, written $\sum A_i$ or
$A_0+A_1+\cdots$, means the set
$\bigl\{\sum_i x_i \mid (x_i)\in\bigoplus A_i\bigr\}$.

The language of sum-sets helps express notions of additive number
theory.  For example, the famous Twin Prime Conjecture predicts an
infinite subset $T\subset\PP$ such that $T+\{0,2\}\subset\PP$.  A long
record of intractability has not prevented the Twin Prime Conjecture
from spawning many (of course conjectural) generalizations, perhaps
most famously Schinzel's Hypothesis H \cite{SchSie}.  Here we develop
a formal consequence of an older generalization of intermediate
ambition, namely Dickson's conjecture \cite{Dickson}.  Dickson's
conjecture concerns finite families $\{a_i m+b_i\}$ of linear
polynomials in $m$ with $a_i\ge 1$ and $b_i\in\NN$,
having the property that no one prime divides $\prod_i(a_i n+b_i)$
for every $n\in\NN$; the conjecture predicts infinitely many values of
$m$ making all the quantities $a_i m+b_i$ simultaneously prime.

Now we can state our results.

\begin{theorem}\label{thm:main}
Assume Dickson's conjecture.  Then for every prime $p$ and every
integer $n$ with $1\le n\le p-2$ there exist infinite sets
$E_0,\dots,E_n\subset 2\NN$, each containing $0$, such that
\[
\{p\}+\sum_{i=0}^{n}E_i \;\subset\; \PP .
\]
\end{theorem}

The restriction $n\le p-2$ cannot be relaxed, and in particular no
infinite-dimensional analogue can hold:

\begin{proposition}\label{prop:sharp}
Let $p$ be prime and let $E_0,\dots,E_{p-1}\subset\NN$ each contain
$0$ and at least one positive element.  Then
$\{p\}+\sum_{i=0}^{p-1}E_i \not\subset \PP$.  Consequently, for no
prime $p$ and no family $\{E_i\}_{i\in I}$ with $|I|\ge p$ (in
particular, no infinite family) of sets each containing $0$ and a
positive element can one have $\{p\}+\sum_{i\in I}E_i\subset\PP$.
\end{proposition}

Together the two statements give a sharp dichotomy: under Dickson's
conjecture, a sum-set of the form $\{p\}+E_0+\cdots+E_n$ (with
$0\in E_i$, $E_i$ infinite) lies within the primes exactly when
$n\le p-2$.

We place these statements in context.  The existence assertion of
Theorem~\ref{thm:main}, viewed simply as an $(n+1)$-fold conditional
sum-set within the primes, is not new: Granville \cite{Granville}
showed already in 1990, under the Dickson--Hardy--Littlewood
conjecture, that the primes contain $k$-fold sum-sets
$A_1+\cdots+A_k$ of infinite sets for every $k$.  Indeed the
unanchored form of Theorem~\ref{thm:main} falls out of Granville's
theorem in a few lines: with $A_0+\cdots+A_n\subseteq\PP$ as
\cite{Granville} provides, choose $a_i\in A_i$ and put
$p' := a_0+\cdots+a_n$, itself a full sum and hence prime; the
translates $E_i := \{\,b-a_i : b\in A_i,\ b\ge a_i,\ b\equiv a_i
\pmod 2\,\}$ are infinite, consist of even numbers, contain $0$, and
satisfy $\{p'\}+E_0+\cdots+E_n \subseteq A_0+\cdots+A_n
\subseteq\PP$.  What this derivation cannot control is the anchor:
$p'$ falls where it falls.  What we
believe is new here is therefore the anchored, quantitative form of the
question and its answer: the exact trade-off between the anchor prime
$p$ and the achievable dimension, with the threshold $n=p-2$ sharp in
both directions.  Tao and Ziegler \cite{TaoZiegler} proved
unconditionally that there exist infinite sets $A,B$ with $a+b$ prime
whenever $a\in A$, $b\in B$ and $a<b$; whether the full sum-set $A+B$
of two infinite sets lies in the primes remains open unconditionally.
Kra, Moreira, Richter and Robertson \cite{KMRRsurvey} prove,
conditionally on Dickson--Hardy--Littlewood, the stronger single-set
statement that $\PP-1$ contains all sums of at most $k$ distinct
elements of a single infinite set, for each fixed $k$; they also show
unconditionally that no shift of $\PP$ contains an IP-set, a statement
in the spirit of Proposition~\ref{prop:sharp}.  We
return in the closing remarks to a second point of interest, the
contrast with random models of the primes.

\section{Proof of Proposition~\ref{prop:sharp}}

Choose a positive element $x_i\in E_i$ for $i=0,\dots,p-1$ and form
the prefix sums $\sigma_k := x_0+\cdots+x_{k-1}$ for $k=0,\dots,p$
(so $\sigma_0=0$).  The $p+1$ values $\sigma_k \bmod p$ cannot all be
distinct, so $\sigma_a\equiv\sigma_b \pmod p$ for some $a<b$, whence
\[
x_a + x_{a+1} + \cdots + x_{b-1} \equiv 0 \pmod p .
\]
Taking the coordinates $a,\dots,b-1$ equal to these $x_i$ and all
other coordinates equal to $0$ produces an element of
$\{p\}+\sum_i E_i$ congruent to $0$ modulo $p$ and strictly larger
than $p$; that element is composite.  The final assertion follows by
restricting attention to any $p$ of the sets.  \qed

\begin{remarks}
The pigeonhole step is the statement that the Davenport constant of
$\mathbb{Z}/p\mathbb{Z}$ equals $p$; the bound is attained (take every
$x_i\equiv 1 \bmod p$ with only $p-1$ sets), which is why the
threshold sits at $p-1$ sets rather than $p$.  Note that the
proposition is unconditional and requires neither infinitude nor
evenness of the $E_i$.
\end{remarks}

\section{Proof of Theorem~\ref{thm:main}}

Fix $n\ge 1$ and a prime $p>n+1$.  We must produce $n+1$ infinite sets
$E_i\subset 2\NN$ of even numbers, each containing $0$.

We notate the elements of $E_i$ from small to large as $e_{i,j}$,
$j=0,1,\dots$, starting with $e_{i,0}=0$ for all $i$, and we determine
the further elements inductively in the order
\[
e_{0,1},\,e_{1,1},\,\dots,\,e_{n,1},\,e_{0,2},\,\dots,\,e_{n,2},\,
e_{0,3},\,\dots
\]
Write $E_{i,j}:=\{e_{i,0},\dots,e_{i,j}\}$ for the provisional
$(j+1)$-term initial segment of $E_i$.  At every stage the
\emph{current sum-set} $S$ means $\{p\}+\sum_i E_{i,j_i}$ formed from
the current initial segments; the construction will maintain
$S\subset\PP$ throughout, and $E_i:=\bigcup_j E_{i,j}$ then completes
the proof.

Roughly, the strategy for determining each $e_{i,j}$ comes down to
applying Dickson's conjecture to translate the set
\[
C^\ast \;:=\; \{p\}+\sum_{i'<i}E_{i',j}+\sum_{i'>i}E_{i',j-1}
\]
into $\PP$ by some even number, which we then declare to be $e_{i,j}$:
the sums newly created by the insertion are exactly
$\{e_{i,j}\}+C^\ast$, since a sum using an older element of coordinate
$i$ was already present.  The point requiring care is that we must
also avoid spoiling the hypothesis of Dickson's conjecture at every
\emph{later} stage of the construction.  We now set up the
bookkeeping that accomplishes this.

\subsection*{Critical values and the residues $s_q$}

Call an element of the current sum-set $S$ a \emph{critical value} if
it admits a representation $p+\sum_{i\in T}x_i$ with
$T\subsetneq\{0,\dots,n\}$ and $x_i\in E_{i,j_i}\setminus\{0\}$ ---
that is, a representation with at least one coordinate equal to $0$.
(If one visualizes the growing sum-set as a box with $p$ at the
origin, the critical values sit on the bounding hyperplanes through
the origin.)  Critical values are primes, being elements of $S$; we
call them \emph{critical primes}.  Once critical, always critical.
The prime $p$ itself (the empty sum) is critical from the outset.
Observe that the set $C^\ast$ above consists entirely of critical
values: its representations omit coordinate $i$.

When a prime $q$ first occurs as a critical value --- its
\emph{birth} --- let $C_q$ denote the set of critical values present
at that moment (including $q$ itself) and put $N_q:=|C_q|$.  Provided
\begin{equation}\label{eq:room}
q \;>\; (n+1)\,N_q ,
\end{equation}
we may choose a residue $s_q\in\{1,\dots,q-1\}$ such that
\begin{equation}\label{eq:sq}
(k+1)\,s_q \;\not\equiv\; -c \pmod q
\qquad\text{for all } 0\le k\le n \text{ and all } c\in C_q :
\end{equation}
each pair $(k,c)$ forbids exactly one residue, since
$k+1\le n+1<p\le q$ is invertible modulo $q$, and there are at most
$(n+1)N_q<q$ such pairs.  (Taking $k=0$ and $c=q$ shows
\eqref{eq:sq} forces $s_q\not\equiv 0$, so the choice
$s_q\in\{1,\dots,q-1\}$ involves no further constraint.)  We treat $p$
as born at the start, with $C_p=\{p\}$ and $N_p=1$; here the forbidden
set is exactly $\{0\}$, so \emph{any} nonzero residue $s_p$ works ---
this is precisely where the hypothesis $p>n+1$ enters.

\emph{Rule for critical primes:} every element $e_{i,j}$ inserted
after the birth of $q$ satisfies $e_{i,j}\equiv s_q \pmod q$.

The utility of \eqref{eq:sq} rests on the following confinement
lemma, which controls all critical values forever after.

\begin{lemma}\label{lem:confine}
Let $q$ be a critical prime.  Every critical value $c'$ arising at any
time satisfies
\[
c' \;\equiv\; c + k\,s_q \pmod q
\qquad\text{for some } c\in C_q \text{ and some } 0\le k\le n .
\]
\end{lemma}

\begin{proof}
Write $c'=p+\sum_{i\in T}x_i$ with $T\subsetneq\{0,\dots,n\}$ and each
$x_i$ positive, and let $k$ be the number of $i\in T$ whose element
$x_i$ was inserted after the birth of $q$.  Each such $x_i$ is
$\equiv s_q\pmod q$ by the rule above, and $k\le|T|\le n$.  Deleting
those coordinates leaves the value $p+\sum x_i$ over the remaining
(pre-birth) coordinates; all of its constituents existed at the birth
of $q$, so it belonged to the sum-set then, and it has an empty
coordinate (any element of $\{0,\dots,n\}\setminus T$, and the deleted
coordinates besides), so it is a critical value present at birth:
an element of $C_q$.
\end{proof}

\subsection*{Frozen primes}

Call a prime $q$ \emph{urgent} at a given step if $q$ is not critical
and $q$ does not exceed the cardinality of the upcoming sum-set.
Urgency is permanent, since the sum-set only grows.

\emph{Rule for urgent primes:} every element inserted from that step
on satisfies $q \mid e_{i,j}$.

Note that $q=2$ is urgent from the very first insertion (the upcoming
sum-set already has two elements, and $2$ is never a critical value,
critical values being odd primes), so all the $e_{i,j}$ are
automatically even and $E_i\subset 2\NN$, as demanded.

\subsection*{The inductive step}

Suppose the construction has run without incident and we are to insert
$e_{i,j}$, with $C^\ast$ as above.  Let $A$ be the product of all
currently urgent primes and all current critical primes, and use the
Chinese Remainder Theorem to find $B$ with $B\equiv 0 \pmod q$ for
each urgent $q$ and $B\equiv s_q\pmod q$ for each critical $q$ (the
moduli are distinct primes).  We seek $e_{i,j}$ of the form $Am+B$;
imposing the congruence rules is then automatic.  Consider the family
of linear forms
\[
\{\,Am+B+c \;:\; c\in C^\ast\,\}.
\]
To apply Dickson's conjecture we must check that no prime $q$ divides
$\prod_{c}(Am+B+c)$ for every $m$.

If $q\nmid A$, then $Am+B$ runs over all residues modulo $q$, so it
suffices that the $c$'s not exhaust the residue classes modulo $q$.
But $q\nmid A$ means $q$ is neither urgent nor critical, whence $q$
exceeds the cardinality of the upcoming sum-set and in particular
$|C^\ast|<q$.

If $q$ is urgent, then $Am+B+c\equiv c\pmod q$, and each $c\in C^\ast$
is a prime different from $q$ (the elements of $C^\ast$ are critical
values and $q$ is not critical), so $q\nmid\prod_c c$.

If $q$ is critical, then $Am+B+c\equiv s_q+c \pmod q$.  By
Lemma~\ref{lem:confine} each $c\in C^\ast$ satisfies
$c\equiv c_0+k s_q$ with $c_0\in C_q$ and $0\le k\le n$, so
$s_q+c\equiv (k+1)s_q+c_0\not\equiv 0\pmod q$ by \eqref{eq:sq}.

Dickson's conjecture therefore supplies infinitely many $m$ making all
$Am+B+c$ simultaneously prime.  Among these we choose $e_{i,j}:=Am+B$
so large that
\begin{itemize}
\item[(i)] every new sum $e_{i,j}+c$ exceeds every element of the old
  sum-set and every currently urgent or critical prime; and
\item[(ii)] every new sum exceeds $(n+1)N+1$, where $N$ is the number
  of critical values after the insertion.
\end{itemize}
Condition (i) guarantees that urgent primes never become critical
(a prime already tracked can never reappear as a new sum) and that the
newly created sums are genuinely new.  Condition (ii) is the
\emph{growing-room} requirement: any prime $q$ born critical at this
insertion is itself one of the new sums, so $q>(n+1)N\ge(n+1)N_q$,
which is exactly \eqref{eq:room}; hence $s_q$ exists and the
bookkeeping can continue.  (Several critical primes may be born at one
insertion; they are then present in one another's sets $C_q$, as
Lemma~\ref{lem:confine} requires.)

All new sums $\{e_{i,j}\}+C^\ast$ are prime, so the enlarged sum-set
again lies within $\PP$, closing the induction.  Since every
coordinate receives infinitely many insertions, each
$E_i=\bigcup_j E_{i,j}$ is infinite, and
$\{p\}+\sum_i E_i=\bigcup S\subset\PP$.  \qed

\section{Remarks}

Both principal results of this paper have been formally verified in
Lean~4 against the Mathlib library by Aristotle (Harmonic), in
collaboration with Claude (Anthropic).
Proposition~\ref{prop:sharp} and its finite-family consequence are
verified unconditionally.  Theorem~\ref{thm:main} is verified in
conditional form: Dickson's conjecture is encoded as an explicit
hypothesis of the formal statement rather than as an axiom, so the
axiom footprint of the entire development consists only of Mathlib's
standard foundational principles, and the unconditional results
cannot silently depend on the conjecture.  The formalization of the
construction proceeds through a symmetric strengthening: an injective
sequence of even numbers all of whose anchored sums of at most $n+1$
distinct terms are prime, in the spirit of the $\mathrm{FS}_{\le
k}$-statement of \cite{KMRRsurvey}, Appendix~B.

The Bateman--Horn conjecture \cite{BatemanHorn} includes a
quantitative version of Dickson's conjecture, but asymptotic densities
without error estimates do not help us bound the $e_{i,j}$.  Lacking
even conjectural tools to bound growth rates, the construction above
aims for clarity rather than parsimony.

Granville's conclusion from Dickson's conjecture is dramatic: the
primes contain $k$-fold sum-sets of infinite sets for every $k$.  One
might well ask whether so strong a conclusion ought to count as
evidence \emph{against} the conjecture.  The prospect has a special
appeal for the skeptic: a contradiction derived from Granville's
conclusion alone would demolish Dickson's conjecture wholesale,
without the labor of singling out a particular fatal admissible
family.  Alas for the skeptic, the probabilistic heuristic supports
Granville's conclusion.  Dickson's conjecture holds in almost every
random model that captures the density behavior of the primes, and
such random sets contain the analogous sum-sets.  But curiously, the
same heuristic supports a conclusion that is provably false.  A
typical random set with prime-like densities contains
\emph{infinite}-dimensional sum-sets: generically one may insert
infinitely many dimension-increasing steps into the construction,
the divisibility obstruction underlying
Proposition~\ref{prop:sharp} having no analogue for a random set.
For the primes themselves, Proposition~\ref{prop:sharp} caps the
dimension of a sum-set anchored at $p$, unconditionally, at $p-2$.
The heuristic thus cannot distinguish the true conclusion from the
false one; what separates them is arithmetic invisible to density
models, namely the congruence obstruction of
Proposition~\ref{prop:sharp}.  The dichotomy of this paper locates
the exact boundary at which the primes part company with their random
models --- and, granting Dickson, the cap $p-2$ is attained.

\begin{thebibliography}{9}

\bibitem{BatemanHorn}
P.~T. Bateman and R.~A. Horn,
\emph{A heuristic asymptotic formula concerning the distribution of
prime numbers},
Math.\ Comp.\ \textbf{16} (1962), 363--367.

\bibitem{Dickson}
L.~E. Dickson,
\emph{A new extension of Dirichlet's theorem on prime numbers},
Messenger of Mathematics \textbf{33} (1904), 155--161.

\bibitem{Granville}
A.~Granville,
\emph{A note on sums of primes},
Canad.\ Math.\ Bull.\ \textbf{33} (1990), no.~4, 452--454.

\bibitem{KMRRsurvey}
B.~Kra, J.~Moreira, F.~K. Richter and D.~Robertson,
\emph{Problems on infinite sumset configurations in the integers and
beyond}, preprint, arXiv:2311.06197.

\bibitem{SchSie}
A.~Schinzel and W.~Sierpi\'nski,
\emph{Sur certaines hypoth\`eses concernant les nombres premiers},
Acta Arith.\ \textbf{4} (1958), 185--208.

\bibitem{TaoZiegler}
T.~Tao and T.~Ziegler,
\emph{Infinite partial sumsets in the primes},
J.\ Anal.\ Math.\ \textbf{151} (2023), 375--389.

\end{thebibliography}

\end{document}
